\documentclass[a4paper,12pt]{article}

\usepackage[T2A]{fontenc}
\usepackage[utf8]{inputenc}
\usepackage[russian]{babel}

\usepackage{amsmath,amssymb,amsthm,mathtools}
\usepackage{helvet}
\renewcommand{\familydefault}{\sfdefault}

\usepackage{icomma}
\usepackage[
  a4paper,
  margin=2.05cm,
  headheight=18pt
]{geometry}
\usepackage{microtype}

\usepackage{tikz}
\usetikzlibrary{calc}

\usepackage{tkz-euclide}

\usepackage{pgfplots}
\pgfplotsset{compat=1.18}

\usepackage{tabularx,booktabs,longtable}
\usepackage{enumitem}
\usepackage{hyperref}
\usepackage{parskip}
\usepackage{fancyhdr}
\usepackage{setspace}
\usepackage{xcolor}

% -------------------------------------------------
% Цветовая палитра
% -------------------------------------------------

\definecolor{accent}{HTML}{008080}
\definecolor{darkaccent}{HTML}{004D4D}
\definecolor{textgray}{HTML}{333333}
\definecolor{softgray}{HTML}{E9EEEE}

% -------------------------------------------------
% Общие настройки
% -------------------------------------------------

\hypersetup{
  colorlinks=true,
  linkcolor=darkaccent,
  urlcolor=darkaccent
}

\onehalfspacing

\setlength{\parindent}{0pt}
\setlength{\emergencystretch}{2em}

\setlist[itemize]{
  leftmargin=1.6em,
  itemsep=0.25em,
  topsep=0.35em
}

\setlist[enumerate]{
  leftmargin=1.8em,
  itemsep=0.4em,
  topsep=0.4em
}

\renewcommand{\arraystretch}{1.28}

\pagestyle{fancy}
\fancyhf{}
\lhead{\small\textcolor{textgray}{София Кальней \textbullet\ ОГЭ}}
\rhead{\small\textcolor{textgray}{Чек-лист по алгебре}}
\cfoot{\small\thepage}

\renewcommand{\headrulewidth}{0.3pt}
\renewcommand{\headrule}{%
  \hbox to\headwidth{%
    \color{softgray}%
    \leaders\hrule height \headrulewidth\hfill
  }%
}

\newcommand{\checkitem}{\item[$\square$]}
\newcommand{\important}[1]{\textcolor{darkaccent}{\textbf{#1}}}
\newcommand{\R}{\mathbb{R}}

\begin{document}

% =================================================
% ТИТУЛЬНЫЙ ЛИСТ
% =================================================

\begin{titlepage}
\thispagestyle{empty}
\centering

\vspace*{2.2cm}

{\Large
\textcolor{darkaccent}{\textbf{Чек-лист}}
\par}

\vspace{0.8cm}

{\huge\bfseries
Рациональные неравенства, радикалы,\\[0.2cm]
текстовые задачи и графики с модулем
\par}

\vspace{0.7cm}

{\large
Подготовка к ОГЭ по математике
\par}

\vfill

{\large
\textbf{Лёвин Артём Александрович}\\[0.15cm]
эксклюзивно для \textbf{Софии Кальней}
\par}

\vspace{0.6cm}

{\large
\today
\par}

\vfill

\begin{tikzpicture}[x=1cm,y=1cm]
\draw[
  accent!55,
  line width=0.9pt
]
(0,0)
.. controls (3.0,0.8) and (6.8,-0.8) .. (9.5,0)
.. controls (11.4,0.5) and (13.0,0.2) .. (15.0,0.55);
\end{tikzpicture}

\end{titlepage}

% =================================================
% 1. МЕТОД ИНТЕРВАЛОВ
% =================================================

\section*{\textcolor{darkaccent}{
1. Метод интервалов: полиномиальные и рациональные неравенства
}}

\subsection*{1.1. Что обязательно проверить}

\begin{itemize}

\checkitem
Разложите числитель и знаменатель на множители настолько,
насколько это возможно.

\checkitem
Отметьте все критические точки на числовой прямой.

\checkitem
Нули числителя при нестрогом неравенстве можно включать в ответ.

\checkitem
Нули знаменателя \important{никогда не входят} в ответ:
в этих точках выражение не определено.

\checkitem
При переходе через корень нечётной кратности знак меняется;
через корень чётной кратности знак сохраняется.

\checkitem
Для калибровки знаков достаточно проверить одну удобную точку
на одном интервале, затем использовать кратности корней.

\end{itemize}

\subsection*{1.2. Пример: повторяющийся корень}

Решим неравенство

\[
(x^2-2x-15)(x^2-7x+10)\le 0.
\]

Разложим квадратные трёхчлены:

\[
x^2-2x-15=(x+3)(x-5),
\]

\[
x^2-7x+10=(x-2)(x-5).
\]

Следовательно,

\[
(x+3)(x-2)(x-5)^2\le0.
\]

Критические точки:

\[
-3,\qquad 2,\qquad 5.
\]

Корень \(x=5\) имеет чётную кратность, поэтому знак
при переходе через него сохраняется.

Для калибровки удобно взять \(x=0\):

\[
(0+3)(0-2)(0-5)^2<0.
\]

Значит, на интервале \((-3;2)\) стоит знак «минус».

\begin{center}
\begin{tikzpicture}[x=1.35cm,y=1cm]

\draw[->,line width=0.7pt]
(-5.2,0)--(6.8,0)
node[right] {$x$};

\foreach \x/\lab in {-3/-3,2/2,5/5} {
  \draw (\x,0.11)--(\x,-0.11);
  \fill[accent] (\x,0) circle (2.2pt);
  \node[below=5pt] at (\x,0) {$\lab$};
}

\node[above=5pt] at (-4.1,0) {$+$};
\node[above=5pt] at (-0.5,0) {$-$};
\node[above=5pt] at (3.5,0) {$+$};
\node[above=5pt] at (5.9,0) {$+$};

\node[
  above=22pt,
  text=darkaccent
] at (5,0)
{кратность \(2\)};

\end{tikzpicture}
\end{center}

Получаем

\[
\boxed{x\in[-3;2]\cup\{5\}}.
\]

Точка \(x=5\) входит в ответ отдельно:
в ней левая часть равна нулю, а неравенство нестрогое.

\subsection*{1.3. Пример: рациональное неравенство}

Рассмотрим

\[
\frac{x^2-3x-4}{(x+10)(x-3)}\ge0.
\]

Разложим числитель:

\[
x^2-3x-4=(x-4)(x+1).
\]

Получаем

\[
\frac{(x-4)(x+1)}
     {(x+10)(x-3)}
\ge0.
\]

Критические точки:

\[
-10,\qquad -1,\qquad 3,\qquad 4.
\]

При этом

\[
x\ne-10,\qquad x\ne3,
\]

поскольку эти значения обращают знаменатель в ноль.

После расстановки знаков:

\[
\boxed{
x\in
(-\infty;-10)
\cup[-1;3)
\cup[4;+\infty)
}.
\]

\subsection*{1.4. Типовые ошибки}

\begin{itemize}

\checkitem
Включать в ответ нуль знаменателя из-за знака
\(\ge\) или \(\le\).

\checkitem
Автоматически менять знак при каждом критическом значении,
не учитывая кратность корня.

\checkitem
Сокращать общий множитель и забывать ограничение,
которое существовало в исходном знаменателе.

\checkitem
Терять отдельный корень чётной кратности
при нестрогом неравенстве.

\checkitem
Ошибаться при первоначальной калибровке знака.

\end{itemize}

% =================================================
% 2. РАДИКАЛЫ
% =================================================

\section*{\textcolor{darkaccent}{
2. Уравнения с одинаковым радикалом в обеих частях
}}

Если одинаковое выражение с квадратным корнем
стоит в обеих частях уравнения, его можно сократить,
\important{предварительно выписав ОДЗ}.

Рассмотрим

\[
2x^2+\sqrt{2-x}
=
3x+14+\sqrt{2-x}.
\]

Сначала определим область допустимых значений:

\[
2-x\ge0.
\]

Отсюда

\[
x\le2.
\]

Теперь одинаковые радикалы можно убрать:

\[
2x^2=3x+14.
\]

Получаем квадратное уравнение

\[
2x^2-3x-14=0.
\]

Вычислим дискриминант:

\[
D=(-3)^2-4\cdot2\cdot(-14)
=9+112
=121.
\]

Тогда

\[
x_{1,2}
=
\frac{3\pm11}{4}.
\]

Следовательно,

\[
x_1=\frac72,
\qquad
x_2=-2.
\]

Проверяем корни по ОДЗ:

\[
\frac72>2,
\]

поэтому первый корень не подходит.

Для второго корня:

\[
-2\le2.
\]

Следовательно,

\[
\boxed{x=-2}.
\]

\important{Главная идея:}
после сокращения одинаковых радикалов ОДЗ сохраняется
и используется как фильтр для найденных корней.

\subsection*{2.1. Универсальный алгоритм}

\begin{enumerate}

\item
Выпишите условие существования каждого квадратного корня.

\item
Упростите уравнение только внутри найденной ОДЗ.

\item
Решите полученное алгебраическое уравнение.

\item
Проверьте каждый найденный корень по ОДЗ.

\item
В ответ запишите только прошедшие проверку значения.

\end{enumerate}

% =================================================
% 3. ДВИЖЕНИЕ
% =================================================

\section*{\textcolor{darkaccent}{
3. Текстовая задача на движение
}}

При составлении таблицы переменную удобно вводить
через величину, которую требуется найти.

\textbf{Пример.}

Автомобиль проехал из \(A\) в \(B\) \(180\) км
с постоянной скоростью.
На обратном пути он увеличил скорость на \(5\) км/ч
и затратил на \(24\) минуты меньше.
Найдите исходную скорость.

Обозначим искомую скорость через \(x\) км/ч.

\begin{table}[h]
\centering
\caption{Перевод условия в математическую модель}

\begin{tabularx}{0.92\linewidth}{@{}lccc@{}}
\toprule
Направление
&
Скорость, км/ч
&
Время, ч
&
Расстояние, км
\\
\midrule

\(A\to B\)
&
\(x\)
&
\(\dfrac{180}{x}\)
&
\(180\)
\\[0.25cm]

\(B\to A\)
&
\(x+5\)
&
\(\dfrac{180}{x+5}\)
&
\(180\)
\\

\bottomrule
\end{tabularx}
\end{table}

Имеем ограничение

\[
x>0.
\]

Переведём минуты в часы:

\[
24\text{ мин}
=
\frac{24}{60}\text{ ч}
=
\frac25\text{ ч}.
\]

Время первого пути больше времени обратного пути на
\(\frac25\) часа:

\[
\frac{180}{x}
-
\frac{180}{x+5}
=
\frac25.
\]

Приведём левую часть к общему знаменателю:

\[
\frac{180(x+5)-180x}{x(x+5)}
=
\frac25.
\]

Отсюда

\[
\frac{900}{x(x+5)}
=
\frac25.
\]

Перемножаем крест-накрест:

\[
4500=2x(x+5).
\]

После деления на \(2\):

\[
2250=x^2+5x.
\]

Получаем

\[
x^2+5x-2250=0.
\]

Дискриминант:

\[
D=5^2-4\cdot1\cdot(-2250)
=25+9000
=9025.
\]

Причём

\[
9025=95^2.
\]

Поэтому

\[
x_{1,2}
=
\frac{-5\pm95}{2}.
\]

Получаем

\[
x_1=45,
\qquad
x_2=-50.
\]

Отрицательная скорость физического смысла не имеет.

Ответ:

\[
\boxed{45\text{ км/ч}}.
\]

\subsection*{3.1. Чек-лист для задач на движение}

\begin{itemize}

\checkitem
В качестве \(x\) по возможности обозначайте именно ту величину,
которую требуется найти.

\checkitem
Используйте формулу

\[
t=\frac{s}{v}.
\]

\checkitem
Все значения времени приведите к одинаковым единицам.

\checkitem
Если движение стало быстрее, время при том же расстоянии уменьшилось.

\checkitem
Учитывайте ограничения:

\[
v>0.
\]

\checkitem
После решения уравнения проверяйте физический смысл корней.

\end{itemize}

% =================================================
% 4. БОЛЬШИЕ КВАДРАТНЫЕ КОРНИ
% =================================================

\section*{\textcolor{darkaccent}{
4. Извлечение корней из больших полных квадратов
}}

\subsection*{4.1. Разбиение на удобные множители}

Если число удаётся представить
в виде произведения полных квадратов, используйте

\[
\sqrt{ab}
=
\sqrt a\,\sqrt b,
\qquad
a\ge0,\quad b\ge0.
\]

Например,

\[
\sqrt{14400}
=
\sqrt{144\cdot100}
=
12\cdot10
=
120.
\]

Этот способ особенно удобен,
когда число оканчивается нулями.

\subsection*{4.2. Факторизация}

Рассмотрим

\[
\sqrt{17424}.
\]

Последовательно раскладываем число:

\[
17424
=
2^4\cdot3^2\cdot11^2.
\]

Поэтому

\[
\sqrt{17424}
=
\sqrt{2^4\cdot3^2\cdot11^2}.
\]

При извлечении квадратного корня
показатели чётных степеней делятся на \(2\):

\[
\sqrt{17424}
=
2^2\cdot3\cdot11.
\]

Значит,

\[
\boxed{\sqrt{17424}=132}.
\]

Ещё один полезный пример:

\[
15376
=
16\cdot961
=
4^2\cdot31^2.
\]

Следовательно,

\[
\sqrt{15376}
=
4\cdot31
=
124.
\]

\subsection*{4.3. Поиск близкого полного квадрата}

Иногда большой дискриминант удобнее распознать
через формулы сокращённого умножения.

Например,

\[
D
=
1005^2-4\cdot3006.
\]

Считаем:

\[
1005^2=1\,010\,025,
\]

\[
4\cdot3006=12\,024.
\]

Поэтому

\[
D
=
1\,010\,025-12\,024
=
998\,001.
\]

Число находится очень близко к

\[
1000^2=1\,000\,000.
\]

Используем формулу

\[
(a-b)^2=a^2-2ab+b^2.
\]

Получаем

\[
999^2
=
(1000-1)^2
=
1\,000\,000-2000+1
=
998\,001.
\]

Следовательно,

\[
\boxed{\sqrt D=999}.
\]

\subsection*{4.4. Последняя цифра как быстрый фильтр}

Полезно помнить:

\[
\begin{array}{c|c}
\text{Последняя цифра числа}
&
\text{Последняя цифра квадрата}
\\
\midrule
0 & 0\\
1 & 1\\
2 & 4\\
3 & 9\\
4 & 6\\
5 & 5\\
6 & 6\\
7 & 9\\
8 & 4\\
9 & 1
\end{array}
\]

Например, если полный квадрат оканчивается цифрой \(1\),
его квадратный корень может оканчиваться цифрой \(1\) или \(9\).

\important{Важно:}
этот признак используется только для отбора кандидатов.
Найденное число необходимо проверить возведением в квадрат.

% =================================================
% 5. МОДУЛЬ И ГРАФИКИ
% =================================================

\section*{\textcolor{darkaccent}{
5. Графики с модулем и параметры
}}

\subsection*{5.1. Раскрытие модуля по определению}

Для выражения \(|x+2|\):

\[
|x+2|
=
\begin{cases}
x+2, & x\ge-2,\\
-(x+2), & x<-2.
\end{cases}
\]

Рассмотрим функцию

\[
y
=
x^2+\frac52x-\frac52|x+2|+1.
\]

При

\[
x\ge-2
\]

получаем

\[
|x+2|=x+2.
\]

Следовательно,

\[
y
=
x^2+\frac52x
-\frac52(x+2)+1.
\]

После упрощения:

\[
y=x^2-4.
\]

При

\[
x<-2
\]

имеем

\[
|x+2|=-(x+2).
\]

Поэтому

\[
y
=
x^2+\frac52x
+\frac52(x+2)+1,
\]

откуда

\[
y=x^2+5x+6.
\]

Итак,

\[
\boxed{
y=
\begin{cases}
x^2-4, & x\ge-2,\\
x^2+5x+6, & x<-2.
\end{cases}
}
\]

\subsection*{5.2. Исследование первой ветви}

Первая ветвь:

\[
y=x^2-4,
\qquad
x\ge-2.
\]

Координата вершины:

\[
x_{\text{в}}
=
-\frac{b}{2a}
=
0.
\]

Тогда

\[
y_{\text{в}}=-4.
\]

Вершина:

\[
V_1(0;-4).
\]

Пересечения с осью \(Ox\):

\[
x^2-4=0.
\]

Отсюда

\[
x=\pm2.
\]

Получаем точки

\[
(-2;0),
\qquad
(2;0).
\]

Обе удовлетворяют условию

\[
x\ge-2.
\]

\subsection*{5.3. Исследование второй ветви}

Вторая ветвь:

\[
y=x^2+5x+6,
\qquad
x<-2.
\]

Координата вершины:

\[
x_{\text{в}}
=
-\frac{5}{2}
=
-\frac52.
\]

Подставим:

\[
y_{\text{в}}
=
\left(-\frac52\right)^2
+
5\left(-\frac52\right)
+
6.
\]

Получаем

\[
y_{\text{в}}
=
\frac{25}{4}
-
\frac{25}{2}
+
6
=
-\frac14.
\]

Вершина:

\[
V_2
\left(
-\frac52;
-\frac14
\right).
\]

Нули функции:

\[
x^2+5x+6=0.
\]

Разложим:

\[
(x+2)(x+3)=0.
\]

Отсюда

\[
x=-2
\quad\text{или}\quad
x=-3.
\]

Значение \(x=-2\) не принадлежит области данной ветви,
поскольку здесь требуется

\[
x<-2.
\]

\begin{center}

\begin{tikzpicture}

\begin{axis}[
  width=0.84\linewidth,
  height=7.2cm,
  axis lines=middle,
  unit vector ratio=1 1 1,
  xmin=-5,
  xmax=4,
  ymin=-5,
  ymax=7,
  samples=220,
  clip=false,
  xtick={-4,-3,-2,-1,0,1,2,3},
  ytick={-4,-2,0,2,4,6},
  xlabel={$x$},
  ylabel={$y$},
  grid=both,
  major grid style={draw=softgray},
  minor tick num=1
]

\addplot[
  accent,
  line width=1.1pt,
  domain=-2:3.2
]
{x^2-4};

\addplot[
  darkaccent,
  line width=1.1pt,
  domain=-4.7:-2.005
]
{x^2+5*x+6};

\addplot[
  only marks,
  mark=*,
  mark size=2.2pt,
  accent
]
coordinates {
  (-2,0)
  (0,-4)
  (2,0)
};

\addplot[
  only marks,
  mark=*,
  mark size=2.2pt,
  darkaccent
]
coordinates {
  (-3,0)
  (-2.5,-0.25)
};

\end{axis}

\end{tikzpicture}

\end{center}

\subsection*{5.4. Горизонтальная прямая \(y=m\)}

После построения графика можно исследовать количество
пересечений с горизонтальной прямой

\[
y=m.
\]

Для данного графика ровно три общие точки появляются
при двух значениях параметра:

\[
\boxed{
m=-\frac14
\quad\text{или}\quad
m=0
}.
\]

При

\[
-\frac14<m<0
\]

горизонтальная прямая имеет уже четыре общие точки с графиком.

\important{Вывод:}
в параметрических задачах необходимо считать все пересечения,
а чертёж следует выполнять достаточно крупно и аккуратно.

% =================================================
% 6. БАЗОВЫЙ V-ОБРАЗНЫЙ ГРАФИК
% =================================================

\section*{\textcolor{darkaccent}{
6. Функция \(y=|x+5|-4\)
}}

Рассмотрим

\[
y=|x+5|-4.
\]

Точка смены формулы определяется условием

\[
x+5=0.
\]

Следовательно,

\[
x=-5.
\]

При

\[
x\ge-5
\]

имеем

\[
y=x+5-4=x+1.
\]

При

\[
x<-5
\]

получаем

\[
y=-(x+5)-4=-x-9.
\]

Итак,

\[
\boxed{
y=
\begin{cases}
x+1, & x\ge-5,\\
-x-9, & x<-5.
\end{cases}
}
\]

Вершина графика:

\[
\boxed{(-5;-4)}.
\]

\begin{center}

\begin{tikzpicture}

\begin{axis}[
  width=0.78\linewidth,
  height=6.5cm,
  axis lines=middle,
  unit vector ratio=1 1 1,
  xmin=-10,
  xmax=4,
  ymin=-6,
  ymax=7,
  samples=220,
  clip=false,
  xtick={-10,-8,-6,-5,-4,-2,0,2,4},
  ytick={-4,-2,0,2,4,6},
  xlabel={$x$},
  ylabel={$y$},
  grid=both,
  major grid style={draw=softgray}
]

\addplot[
  accent,
  line width=1.15pt,
  domain=-5:3.5
]
{x+1};

\addplot[
  darkaccent,
  line width=1.15pt,
  domain=-9.5:-5
]
{-x-9};

\addplot[
  only marks,
  mark=*,
  mark size=2.4pt,
  accent
]
coordinates {
  (-5,-4)
};

\end{axis}

\end{tikzpicture}

\end{center}

\subsection*{6.1. Пересечения с \(y=m\)}

Для горизонтальной прямой получаем:

\[
\begin{array}{c|c}
\text{Число общих точек}
&
\text{Условие}
\\
\hline
0 & m<-4\\
1 & m=-4\\
2 & m>-4
\end{array}
\]

% =================================================
% 7. ПАРАМЕТР y=kx
% =================================================

\section*{\textcolor{darkaccent}{
7. Прямая \(y=kx\) и число общих точек
}}

Прямая

\[
y=kx
\]

всегда проходит через начало координат,
а коэффициент \(k\) задаёт её наклон.

Для корректного определения числа пересечений с графиком

\[
y=|x+5|-4
\]

следует рассматривать обе ветви отдельно.

\subsection*{7.1. Правая ветвь}

Правая ветвь имеет вид

\[
y=x+1,
\qquad x\ge-5.
\]

Точки пересечения с \(y=kx\) удовлетворяют

\[
kx=x+1.
\]

Следовательно,

\[
(k-1)x=1.
\]

При

\[
k\ne1
\]

получаем

\[
x=\frac{1}{k-1}.
\]

Однако найденная точка принадлежит правой ветви
только при условии

\[
\frac{1}{k-1}\ge-5.
\]

Из этого условия следует:

\[
k\le\frac45
\quad\text{или}\quad
k>1.
\]

\subsection*{7.2. Левая ветвь}

Левая ветвь:

\[
y=-x-9,
\qquad x<-5.
\]

Пересечение с \(y=kx\):

\[
kx=-x-9.
\]

Следовательно,

\[
(k+1)x=-9.
\]

При

\[
k\ne-1
\]

получаем

\[
x=-\frac{9}{k+1}.
\]

Необходимо также условие

\[
-\frac{9}{k+1}<-5.
\]

Оно выполняется при

\[
-1<k<\frac45.
\]

\subsection*{7.3. Итоговая классификация}

После учёта обеих ветвей:

\[
\boxed{
\begin{aligned}
0\text{ общих точек:}\quad
&
\frac45<k\le1,
\\[0.2cm]
1\text{ общая точка:}\quad
&
k\le-1
\quad\text{или}\quad
k=\frac45
\quad\text{или}\quad
k>1,
\\[0.2cm]
2\text{ общие точки:}\quad
&
-1<k<\frac45.
\end{aligned}
}
\]

\important{Ключевая проверка:}
если прямая параллельна одной ветви,
это ещё не определяет число пересечений со всем графиком.
Всегда требуется отдельно проверить вторую ветвь.

% =================================================
% 8. ОБЩИЙ АЛГОРИТМ
% =================================================

\section*{\textcolor{darkaccent}{
8. Универсальный алгоритм для графиков с модулем
}}

\begin{enumerate}

\item
Найдите значение \(x\), при котором подмодульное выражение
равно нулю.

\item
Раскройте модуль отдельно на каждом промежутке.

\item
Получите кусочно заданную функцию.

\item
Для каждой параболы найдите:

\begin{itemize}
\item вершину;
\item точки пересечения с \(Ox\);
\item при необходимости точку пересечения с \(Oy\).
\end{itemize}

\item
Для прямой достаточно двух точек.

\item
Обязательно учитывайте область существования каждой ветви.

\item
При построении отмечайте масштаб и подписывайте оси.

\item
В задаче с параметром определите критические положения прямой.

\item
Между критическими положениями проверяйте количество общих точек.

\item
Граничные значения параметра рассматривайте отдельно.

\end{enumerate}

% =================================================
% ТРЕНИРОВОЧНЫЙ БЛОК
% =================================================

\newpage

\section*{\textcolor{darkaccent}{
Тренировочный блок --- Путь А
}}

\textbf{Средний уровень}

\begin{enumerate}[label=\textbf{\arabic*.}]

\item
Решите неравенство

\[
(x^2+x-6)(x^2-7x+12)\le0.
\]

\item
Решите неравенство

\[
\frac{x^2-x-6}{(x+4)(x-1)}\ge0.
\]

\item
Решите уравнение

\[
x^2+\sqrt{3-x}
=
5x-4+\sqrt{3-x}.
\]

\end{enumerate}

\textbf{Выше среднего}

\begin{enumerate}[label=\textbf{\arabic*.},start=4]

\item
Автомобиль проехал \(150\) км с постоянной скоростью.
На обратном пути он увеличил скорость на \(10\) км/ч
и затратил на дорогу на \(30\) минут меньше.
Найдите исходную скорость автомобиля.

\item
Извлеките квадратный корень,
предварительно разложив число на множители:

\[
\sqrt{20736}.
\]

\item
Вычислите без калькулятора,
используя близкий полный квадрат:

\[
\sqrt{3\,229\,209}.
\]

\item
Для функции

\[
y=|x-3|-2
\]

укажите значения \(m\), при которых прямая \(y=m\)
имеет с графиком:

\begin{enumerate}[label=\arabic*)]
\item \(0\) общих точек;
\item \(1\) общую точку;
\item \(2\) общие точки.
\end{enumerate}

\item
Дана функция

\[
f(x)
=
x^2+\frac32x-\frac32|x+1|.
\]

Найдите все значения \(m\),
при которых прямая

\[
y=m
\]

имеет с графиком функции ровно три общие точки.

\item
Для функции

\[
y=|x+4|-3
\]

определите все значения \(k\),
при которых прямая

\[
y=kx
\]

имеет с графиком:

\begin{enumerate}[label=\arabic*)]
\item \(0\) общих точек;
\item \(1\) общую точку;
\item \(2\) общие точки.
\end{enumerate}

\end{enumerate}

\textbf{Сложный уровень}

\begin{enumerate}[label=\textbf{\arabic*.},start=10]

\item
Решите неравенство

\[
\frac{(x-2)^2(x+3)}
     {(x+1)(x-4)^2}
\le0.
\]

Обязательно учтите кратности корней
и точки, в которых знаменатель равен нулю.

\end{enumerate}

% =================================================
% ОТВЕТЫ
% =================================================

\newpage

\section*{\textcolor{darkaccent}{
Ответы к тренировочному блоку
}}

\begin{table}[h]
\centering
\caption{Краткие ответы}

\begin{tabularx}{\linewidth}{@{}cX@{}}

\toprule
№ & Ответ \\
\midrule

1 &
\(
[-3;2]\cup[3;4]
\).
\\

2 &
\(
(-\infty;-4)
\cup[-2;1)
\cup[3;+\infty)
\).
\\

3 &
\(
x=1
\).
\\

4 &
\(50\) км/ч.
\\

5 &
\(144\).
\\

6 &
\(1797\).
\\

7 &
1) \(m<-2\);
2) \(m=-2\);
3) \(m>-2\).
\\

8 &
\(
m=-\dfrac34
\)
или
\(
m=-\dfrac12
\).
\\

9 &
1)
\(
\dfrac34<k\le1
\);

2)
\(
k\le-1
\),
или
\(
k=\dfrac34
\),
или
\(
k>1
\);

3)
\(
-1<k<\dfrac34
\).
\\

10 &
\(
[-3;-1)\cup\{2\}
\).
\\

\bottomrule

\end{tabularx}
\end{table}

\end{document}